%----------------------------------------------------------------------------------------
% General-purpose CV — COMPUTATIONAL BIOLOGY track (quick applications)
%----------------------------------------------------------------------------------------

\documentclass[10pt,a4paper,sans]{moderncv}

\usepackage[utf8]{inputenc}
\usepackage[T1]{fontenc}
\usepackage{lmodern}

\moderncvstyle{classic}
\moderncvcolor{blue}

\usepackage[scale=0.78]{geometry}
\usepackage{ragged2e}

%----------------------------------------------------------------------------------------

\firstname{Kundai Farai}
\familyname{Sachikonye}

\address{Zeltnerstrasse 30}{90443, Nuernberg, Germany}
\mobile{(+49) 1626386344}
\email{kundai.sachikonye@bitspark.com}
\homepage{github.com/fullscreen-triangle}
\photo[70pt][0.4pt]{pictures/Bew2.jpg}

\quote{Computational biologist: multi-omics, systems biology, and mass spectrometry,
with machine learning and strong software engineering --- from method to deployed tool.}

%----------------------------------------------------------------------------------------

\begin{document}

\makecvtitle

%----------------------------------------------------------------------------------------

\section{Profile}
\cvitem{}{
  Computational biologist working across genomics / transcriptomics, proteomics /
  metabolomics, and systems-biology modelling. I combine machine learning with
  rigorous software engineering, and validate against authoritative reference data.
}

\section{Technical Skills}
\cvitem{Omics}{
  Genomics, transcriptomics, proteomics, metabolomics; multi-omics integration
}
\cvitem{Systems biology}{
  Cellular and circuit modelling; pathway and network analysis
}
\cvitem{ML \& statistics}{
  PyTorch, scikit-learn; statistical modelling; representation learning
}
\cvitem{Software \& data}{
  Python (primary), R, Rust, SQL; HPC (Ray, Dask, SLURM); reproducible pipelines;
  database design (lipidomics)
}
\cvitem{Validation}{
  NIST, IEDB, 1000~Genomes / gnomAD / UK~Biobank
}

\section{Experience}
\cventry{2021--present}{Independent Computational Biology Researcher}{Bitspark GmbH}{}{}{
  Multi-omics, systems-biology, and mass-spectrometry tools; deployed, documented,
  reference-validated.
}
\cventry{2019--2021}{Computational Lipidomics}{Technische Universität München}{Munich}{}{
  High-throughput mass spectrometry and multi-omics analysis at HPC scale. Taught
  at Master's level.
}
\cventry{2017--2018}{Clinical Glycomics \& Proteomics}{Universitätsklinikum Hamburg-Eppendorf}{Hamburg}{}{
  LC-MS/MS and MALDI-TOF analysis for clinical samples.
}
\cventry{2013}{Cancer Biomarker Discovery}{University Medical Center Groningen}{Groningen}{}{
  UPLC-MS/MS shotgun proteomics on tumour samples.
}

\section{Selected Projects}
\cvitem{gospel}{
  Multi-omics integration (genotype + transcriptomics + metagenomics); validated
  against 1000~Genomes, gnomAD, UK~Biobank.
  \href{https://github.com/fullscreen-triangle/gospel}{github.com/fullscreen-triangle/gospel}
}
\cvitem{systems-biology-shaders}{
  Computational cell / systems-biology modelling (GPU pipeline).
  \href{https://systems-biology-shaders.vercel.app/}{systems-biology-shaders.vercel.app}
}
\cvitem{syndrome}{
  Computational oncology --- neoantigen / MHC prediction (AUC~0.81, IEDB).
  \href{https://syndrome-seven.vercel.app/tools}{syndrome-seven.vercel.app}
}
\cvitem{Crown Prince}{
  Computational pharmacology (PK / PD / ADME); NIST-validated.
  \href{https://crown-prince-four-courners.vercel.app/}{crown-prince-four-courners.vercel.app}
}

\section{Education}
\cventry{2018--2019}{MSc Computational Biology}{Universität Tübingen}{}{}{}
\cventry{2014--2017}{MSc Pharmaceutical Biotechnology}{HAW Hamburg}{}{}{}
\cventry{2011--2014}{BSc Biochemistry and Cell Biology}{Jacobs University Bremen}{}{}{}
\cvitem{}{\textit{Additional:} doctoral research, computational lipidomics, TUM (not completed).}

\section{Languages}
\cvitemwithcomment{English}{Native / Fluent (C2)}{}
\cvitemwithcomment{German}{Conversational}{resident in Germany since 2011}

\renewcommand{\listitemsymbol}{-~}

\end{document}
